\paragraph{How we checked nonconvexity.}
A convex training loss at the midpoint of any two weight vectors is at most the average of their losses. We check this inequality on the actual training examples. A linear recurrent network is linear in its input for fixed weights, but its input and readout weights multiply each other. For example, switch off the state transition and use just one state coordinate. Let its input weight be \(a\) and its readout weight be \(b\). Each tick then outputs \(ab\) times the current input. The pairs \((a,b)=(1,4)\) and \((4,1)\) have the same product, 4, while their midpoint \((5/2,5/2)\) has product \(25/4\). Reversing both readout signs if needed makes increasing the magnitude of this product increase the squared training loss. To see why, the initial slope of the loss as a function of the product at zero can have either sign. Choose the direction with nonnegative slope. Along that direction the loss is a strictly convex quadratic because at least one observed input is nonzero. The midpoint therefore has strictly greater loss than either endpoint.

The same construction works with the bias and readout of one active scalar rectifier, the input and output weights of one sorting block, or the initial state and readout of an autonomous sequence generator. All unused paths can be set to zero. In the sorter the original input still passes through its fixed identity connection. For the two autonomous tasks, a transition of one on the active coordinate keeps the output constant across the observed indices. Exact rational evaluation of the stored training data certifies 36 smooth-function settings (four tasks, widths 1, 4 and 16), six running sums (widths 1 and 2), nine four-input sorters (four, five or six blocks), and 18 sequence generators (Fibonacci widths 1--3 and triangular widths 2--4). Each size uses three starts. The sequence generators share their training labels across starts. Replaying the training code gives midpoint violations from 0.872 to 16.014 in units of training loss. The original 18 affine arithmetic settings and three memoryless sums minimize a weighted sum of squares of affine functions of their weights, a convex quadratic. Addition and subtraction on trainable rectifier networks of widths 1, 2 and 4 are a separate 18-setting experiment with their own midpoint certificates. A midpoint violation establishes the shape of the objective. Whether a solver reaches a better fit is a separate measurement.
